
\appendix
\chapter{Appendix}

This appendix provides key implementation snippets from the CyberSentinel system, including machine learning training, feature extraction, and API integration.

\section{URL Feature Extraction Code}

\begin{lstlisting}[language=Python, caption=URL Feature Extraction for CyberSentinel, label=lst:feature_extraction]
	import re
	from urllib.parse import urlparse
	
	def extract_features(url):
	features = {}
	
	features['url_length'] = len(url)
	features['num_special_chars'] = len(re.findall(r'[@\-_\.]', url))
	features['has_https'] = 1 if url.startswith("https") else 0
	
	parsed = urlparse(url)
	features['num_subdomains'] = len(parsed.netloc.split('.')) - 2
	
	suspicious_keywords = ['login', 'secure', 'verify', 'bank']
	features['has_suspicious_keywords'] = int(
	any(word in url.lower() for word in suspicious_keywords)
	)
	
	return features
\end{lstlisting}

\section{Machine Learning Training Code}

\begin{lstlisting}[language=Python, caption=Random Forest Model Training for URL Classification, label=lst:ml_training]
	import pandas as pd
	from sklearn.ensemble import RandomForestClassifier
	from sklearn.model_selection import train_test_split
	import joblib
	
	def load_dataset(path):
	data = pd.read_csv(path)
	X = data.drop("label", axis=1)
	y = data["label"]
	return X, y
	
	def train_model(X, y):
	X_train, X_test, y_train, y_test = train_test_split(
	X, y, test_size=0.2, random_state=42
	)
	
	model = RandomForestClassifier(
	n_estimators=100,
	random_state=42
	)
	
	model.fit(X_train, y_train)
	
	joblib.dump(model, "phishing_model.pkl")
	return model
\end{lstlisting}

\section{Model Retraining Pipeline}

\begin{lstlisting}[language=Python, caption=Feedback-based Model Retraining Pipeline, label=lst:retrain_pipeline]
	def retrain_model():
	base_data = pd.read_csv("data/url_dataset.csv")
	feedback_data = pd.read_csv("data/approved_feedback.csv")
	
	combined_data = pd.concat([base_data, feedback_data])
	
	X = combined_data.drop("label", axis=1)
	y = combined_data["label"]
	
	model = train_model(X, y)
	print("Model retrained successfully")
	
	return model
\end{lstlisting}

\section{Flask API Endpoint Example}

\begin{lstlisting}[language=Python, caption=Flask API Endpoint for URL Analysis, label=lst:flask_api]
	from flask import Flask, request, jsonify
	
	app = Flask(__name__)
	
	@app.route("/api/analyze-url", methods=["POST"])
	def analyze_url():
	data = request.get_json()
	url = data.get("url")
	features = extract_features(url)
	prediction = model.predict([list(features.values())])[0]
	return jsonify({
		"url": url,
		"prediction": int(prediction)
	})
\end{lstlisting}

\section{Admin Model Retraining Endpoint}

\begin{lstlisting}[language=Python, caption=Admin Endpoint for Model Retraining, label=lst:admin_retrain]
	@app.route("/api/admin/retrain-model", methods=["POST"])
	def retrain():
	try:
	retrain_model()
	return jsonify({"message": "Model retrained successfully"})
	except Exception as e:
	return jsonify({"error": str(e)}), 500
\end{lstlisting}
